身份仲裁与完成记录已经连成一条处理路径。现在沿这条路径逐处中断：写入身份之后退出，扣款之后退出，数据库提交以后、响应以前退出，是否都只会留下允许的业务状态？

一次成功调用只经过一条执行历史。要检查任意次允许的中断，需要列出状态与执行规则，说明每一步怎样保持同一组关系，再用实际中断与有限枚举寻找遗漏。下面继续使用原完整身份 S、输入与结果版本，各个演算就地给出独立初始数据。

\textbf{本章范围为单一 PostgreSQL 18 数据库中的购买与本地扣款，采用 Read Committed；恢复时 E1 有效、身份记录可读、已确认提交持久保留，维护程序不改写终态。}参与并发的资金动作只有购买扣款，这项范围将用于确定余额不足的生效位置。充值扩展、会话关闭与结果回收分别需要补充允许迁移及其保持论证。
